% =============================================================================
% cap4 §4.5 Reprodutibilidade (compactada; códigos movidos p/ Apêndice B)
% =============================================================================
\section{Reprodutibilidade}
\label{sec:metodologia-reprodutibilidade}

O ambiente de execução adotado é Python em ambiente virtual isolado, com arquivo de \emph{lock} de dependências (\texttt{requirements.txt} ou \texttt{environment.yml}) versionado junto ao repositório. Os pesos pré-treinados são baixados explicitamente para diretório local antes da execução, com \emph{checksums} registrados em \texttt{checksums.txt}, permitindo verificar que o modelo carregado em qualquer reexecução é idêntico ao da rodada original. Quatro grupos de dependências precisam ser fixados nominalmente: a \emph{toolkit} de \emph{deep learning} (PyTorch, com \emph{tag} de \emph{build} CUDA quando aplicável), as bibliotecas dos modelos (\modeloMD, \modeloSN, \modeloPetFace, \modeloWildlife), as de manipulação de imagem (OpenCV, Pillow) e as de avaliação (scikit-learn, \emph{pycocotools}). A fixação de \emph{seed} ocorre em três pontos --- inicialização de geradores pseudoaleatórios do Python e do NumPy, \emph{seed} da \emph{toolkit} de \emph{deep learning} e embaralhamento de partições --- a partir de um valor base único registrado no arquivo de configuração; experimentos com múltiplas rodadas (não previstos nesta etapa, mas tratados como trabalho futuro) usam \emph{seeds} consecutivos.

A organização operacional prevê \textbf{cinco scripts} parametrizados por um único arquivo de configuração (JSON ou YAML) versionado junto ao repositório. O \texttt{run\_phase2.py} carrega os \emph{splits} de teste oficiais de CCT e Snapshot Serengeti, executa \modeloMD{} e o baseline YOLO genérico, aplica o mascaramento por \emph{blur} das \emph{bounding boxes} da classe ``pessoa'' (LGPD, Seção~\ref{sec:pipeline-fase2}) antes de qualquer persistência, e produz em \texttt{out/phase2/} um CSV por modelo e estrato com mAP estratificada e custo médio. O \texttt{run\_phase3.py} carrega os subconjuntos relevantes de CCT e NACTI, executa \modeloSN{} e o baseline binário, e produz $F_1$, acurácia balanceada e matriz de confusão. O \texttt{run\_phase4.py} carrega os \emph{splits} oficiais de PetFace e WildlifeReID-10k, monta os catálogos \emph{closed-set} e \emph{open-set} e executa as configurações comparadas (face, flanco por partes do corpo, fusão tardia, baseline global). O \texttt{run\_end2end.py} monta a janela sintética, executa pipeline em cascata e variante ingênua e gera os dados da curva de eficiência. Finalmente, o \texttt{compare\_runs.py} consolida duas execuções em uma única tabela comparável, sinalizando regressões acima de um limiar configurável. Cada execução emite, ao final, um arquivo \texttt{run\_metadata.json} contendo SHA do código, versão de Python, plataforma, \emph{seed} e \emph{hash} da configuração. Os códigos completos dos cinco scripts estão no Apêndice~\ref{ap:pseudocodigos}.
